% Propietario conservado: /Users/ruben/Documents/New project/output/EXTREMA_MEDIA_RAZON_RECIPROCIDAD_ES_EN_20260918/ARTICULO/ES/source/sections/registro_k.tex
\section{Extracción dodecafásica de la constante de acoplamiento aritmético-geométrico}
\label{x_reciprocidad:sec:k-registro}

\begin{definition}[Constante holográfica de acoplamiento aritmético-geométrico]
\label{x_reciprocidad:lib01:definicion}
Se denomina así al registro finito ordenado producido por la extracción
orientada del estado APP--TRIT--TPK, con su origen de fase, normalización y
lectores declarados, cuya representación posicional interviene en la
clausura aritmética y cuyas realizaciones incidencial y geométrica actúan
sobre el mismo registro. El símbolo \(K\) designa su representación
integral canónica; \(\kappa_{\rm ag}\) designa su lectura racional
periódica en la carta fijada.
\end{definition}

La generación de las coordenadas regionales no agota la información del
estado APP--TRIT--TPK. Una lectura arquimediana determina un valor mediante
cilindros compatibles; otra lectura conserva la oposición de las hojas, la
orientación de la ruta y el registro de las incidencias atravesadas. El
registro dodecafásico pertenece a esta segunda lectura antes de intervenir en
la clausura de alfa. Por ello no se define restando una constante física a
una combinación de valores previamente conocidos. Se reconstruye desde una
coordenada firmada del mismo estado aritmético con memoria de trayectoria que produce las coordenadas
regionales.

La distinción es necesaria incluso cuando dos objetos tienen doce
coordenadas. Una palabra de doce bloques del catálogo, una lista de doce
eventos orientados, un vector entero firmado y un registro dodecafásico en base mil son
objetos distintos. Su comparación requiere las aplicaciones que los
relacionan. La igualdad de longitud no proporciona esas aplicaciones, y la
pérdida de signos no es una mera simplificación de notación.

Este capítulo mantiene el corte causal fijado en los capítulos precedentes:
las nueve posiciones APP de los ejes horizontal y vertical producen las
hojas aritméticas; TRIT tipa régimen y orientación; TPK selecciona,
transporta, actualiza y prolonga, conservando residuo, cociente, acarreo,
hoja, ruta, frontera y memoria. Desde la estructura discreta conjunta del
continuo se obtienen las coordenadas que se utilizan aquí. La cadena
\(w_6\to w_{12}\to w_{18}\to w_{24}\to w_{30}\to R_{36}\to G_9\)
no se sustituye por un calendario finito. Tampoco se invierte el orden
\(U_t\to\Gamma_9\to K_{\mathrm{ph}}\): la última aplicación es una
lectura reducida de memoria, no el generador del registro firmado.


% BEGIN OWNER /Users/ruben/Documents/New project/output/EXTREMA_MEDIA_RAZON_RECIPROCIDAD_ES_EN_20260918/ARTICULO/ES/source/sections/revision_k.tex
% Adición propuesta al inicio de registro_k.tex, después de su introducción.
% Conserva todas las definiciones, fórmulas y demostraciones actuales.
% Procedencia: registro_k.tex, §§k-hadamard/k-transversal/k-kappas;
% k_reversibilidad.tex; incidencia_articulacion.tex, inc-ampl-cuaterna.
\paragraph{Coordenatización reversible y determinación incidencial.}
\label{x_reciprocidad:ampl-alfa-k-intro-reversible}
El resultado que organiza este capítulo es la equivalencia entre tres
coordenatizaciones de un registro previamente producido por
APP--TRIT--TPK. La transformación de Hadamard relaciona el registro
firmado con el registro entero; el extractor de diferencias relaciona
este último con sus variaciones transversales y su componente uniforme:
\[
 U\underset{\mathcal H_{12}}{\overset{\mathcal H_{12}/4}{\longleftrightarrow}}K
 \overset{(D_3,D_4,Q)}{\longleftrightarrow}
 (D_3K,D_4K,Q(K)).
\]
La primera equivalencia se establece sobre la subred integral
$\mathcal L_H$; la segunda, sobre la imagen del extractor. Fijadas la
base $1000$, las $12$ posiciones y el origen de fase, la evaluación
$K\mapsto\kappa_{\mathrm{per}}$ añade una cuarta coordenatización
reversible: multiplicar por $1000^{12}-1$ y efectuar las divisiones
euclídeas recupera el registro completo. Las demostraciones de
\S\ref{x_reciprocidad:subsec:k-hadamard}, \S\ref{x_reciprocidad:subsec:k-transversal} y
\S\ref{x_reciprocidad:subsec:k-kappas} establecen estas equivalencias en sus respectivos
dominios.

Sobre ese registro actúan después dos composiciones. La composición
posicional combina sus componentes con las coordenadas de clausura,
propagación y autoescala, y su normalización determina $\alpha_A$. La
composición incidencial aplica el umbral cúbico y obtiene
\[
 K\longmapsto B_K=\{j:K_j\geq729\}=\{4,6,9,10\}.
\]
El registro íntegro permite reconstruir ambas composiciones. El soporte
$B_K$ conserva, específicamente, la pertenencia a la clase seleccionada
por el umbral. La identidad $B_K=H_e\cap P_\pi$ de
\eqref{x_reciprocidad:inc-ampl-cuaterna} enlaza esta extracción entera con la incidencia
de los códigos regionales. El soporte negativo de su balance con $K$
completa la bandera $B_K\subset H^-_\alpha$ que se desarrolla en la
prolongación excepcional. Así, la función de $K$ queda definida mediante
una transformación invertible y mediante aplicaciones posteriores
precisas; la información adicional de profundidad y transporte permanece
en el estado del que procede el registro.

% END OWNER


\subsection{Extensión cíclica y graduación de memoria}
\label{x_reciprocidad:subsec:k-calendario}

El calendario observable contiene ciento ocho posiciones, organizadas en
doce ventanas nonádicas. Con índices positivos, sean
\[
 E_m=\{9(m-1)+1,\ldots,9m\},\qquad 1\leq m\leq12.
\]
El soporte ordenado \(E_1\sqcup\cdots\sqcup E_{12}\) conserva posición y
pertenencia a ventana. El grupo \(C_{108}=\mathbb Z/108\mathbb Z\)
conserva, en cambio, la ley de traslación cíclica. Para el representante
\(0\leq\widetilde\theta<108\), las coordenadas
\[
 \beta(\theta)=\left[\left\lfloor\frac{\widetilde\theta}{9}\right\rfloor\right]_{12},
 \qquad r(\theta)=1+(\widetilde\theta\bmod9)
\]
separan el sector dodecafásico y la posición nonádica interior. Ninguna de
ellas conserva por sí sola el número total de vueltas.

\begin{proposition}[El acarreo del calendario]
\label{x_reciprocidad:prop:k-extension}
Las aplicaciones
\[
 \iota:C_{12}\longrightarrow C_{108},\quad[b]\longmapsto[9b],
 \qquad p:C_{108}\longrightarrow C_9,\quad[\theta]\longmapsto[\theta]_9
\]
forman una sucesión exacta no escindida
\[
 0\longrightarrow C_{12}\xrightarrow{\iota}C_{108}
 \xrightarrow{p}C_9\longrightarrow0.
\]
Para la sección conjuntista que elige los representantes
\(0,\ldots,8\), el defecto de aditividad es el acarreo
\[
 c(r,r')=\left[\left\lfloor\frac{\widetilde r+\widetilde r'}9\right\rfloor\right]_{12}.
\]
\end{proposition}

\begin{proof}
El núcleo de \(p\) está formado por los múltiplos de nueve y coincide con
la imagen de \(\iota\). La reducción es sobreyectiva y la multiplicación
por nueve es inyectiva en el dominio indicado. Si la sucesión se
escindiera, el grupo central sería isomorfo a \(C_{12}\times C_9\), cuyo
exponente es \(\operatorname{mcm}(12,9)=36\); el exponente de
\(C_{108}\) es 108. La contradicción prueba la no escisión. Finalmente,
la división euclídea de \(\widetilde r+\widetilde r'\) por nueve da el
acarreo escrito y verifica la identidad de la sección.
\end{proof}

El resultado explica una diferencia que reaparecerá en alfa: regresar a la
fase no equivale a regresar al estado. Si \(q_{\mathrm{mem}}\) cuenta las
vueltas nonádicas, el calendario conserva solamente
\(\beta=[q_{\mathrm{mem}}]_{12}\). Después de ciento ocho pasos,
\(\beta\) retorna, pero \(q_{\mathrm{mem}}\) aumenta en doce. El
registro conserva además el cilindro, el registro de incidencias y la memoria
vertical. El cierre de una carta finita no convierte esa historia en un
ciclo sin memoria.

\subsection{El registro orientado y la carta integral \texorpdfstring{\(1+5+6\)}{1+5+6}}
\label{x_reciprocidad:subsec:k-registro-integral}

Para una historia producida por TPK, sea \(\tau_m\) la subruta sobre
\(E_m\), \(\sigma_m\) su orientación, \(\kappa_m\) el acarreo de
frontera y \(\Lambda_m\) el artículo local de incidencias. El registro es
\[
 \mathcal R_{12}(x)=
 \bigl((E_m,\tau_m,\sigma_m,\kappa_m,\Lambda_m)\bigr)_{m=1}^{12}.
\]
Su dominio es el espacio de historias aritméticas con memoria de trayectoria, no el conjunto de
etiquetas de fase. La proyección al calendario olvida precisamente algunos
de los datos que el lector firmado necesita.

Una carta inicial del registro puede escribirse mediante eventos
orientados. La extracción orientada fija el conjunto de códigos de evento
\(\{2,4,6,8\}\) y la orientación sectorial
\[
 \Sigma_{12}=(+,+,+,-,-,-,+,+,+,-,-,-).
\]
Si \(d_t\) es el código de dirección en el paso \(t\), se obtiene
\[
 e_i=\Sigma_{12}(i+1)
 \#\{t\in E_{i+1}:d_t\in\{2,4,6,8\}\},
 \qquad0\leq i<12.
\]
La fórmula mantiene separados el censo no negativo de eventos y la
orientación. No identifica una dirección cardinal con un signo TRIT. El
vector \(e\), cuyos componentes son censos locales firmados, tampoco es
todavía el vector \(U\) que aparecerá más abajo.

Las posiciones opuestas determinan, para \(0\leq i<6\),
\[
 p_i=e_i+e_{i+6},\qquad a_i=e_i-e_{i+6}.
\]
Con
\[
 s_C=\sum_{i=0}^{5}p_i,\qquad M_i=p_i-p_5\quad(0\leq i<5),
\]
se define la carta integral
\[
 \mathcal L_{12}(e)
 =(s_C,M_0,\ldots,M_4,a_0,\ldots,a_5).
\]
La descomposición \(1+5+6\) no es un reparto de coordenadas arbitrario:
una coordenada conserva la carga, cinco comparan los pares respecto de un
par de referencia y seis conservan la oposición de las semicoronas.

\begin{lemma}[Inversión de la carta integral]
\label{x_reciprocidad:lem:k-carta-integral}
Una tupla entera \((s_C,M_0,\ldots,M_4,a_0,\ldots,a_5)\) pertenece a la
imagen de \(\mathcal L_{12}\) si y sólo si
\[
 s_C-\sum_{i=0}^{4}M_i\equiv0\pmod6,
 \qquad p_i\equiv a_i\pmod2\quad(0\leq i<6),
\]
donde
\[
 p_5=\frac{s_C-\sum_{i=0}^{4}M_i}{6},\qquad p_i=p_5+M_i\quad(i<5).
\]
En ese caso su preimagen es única y viene dada por
\[
 e_i=\frac{p_i+a_i}{2},\qquad e_{i+6}=\frac{p_i-a_i}{2}.
\]
\end{lemma}

\begin{proof}
La suma de los cinco márgenes es \(s_C-6p_5\), de donde se obtiene la
división por seis. Reconstruidos los \(p_i\), cada par de ecuaciones
\(p_i=e_i+e_{i+6}\), \(a_i=e_i-e_{i+6}\) tiene la solución indicada.
Esta solución es entera exactamente bajo la congruencia de paridad. Todas
las operaciones se invierten, por lo que no queda libertad residual.
\end{proof}

Esta reversibilidad muestra qué información se conserva; no autoriza a
suprimir el resto del artículo. El registro completo tiene un dominio mayor
que esta carta de censos. La extracción de los canales que alimentan a
\(U\) utiliza ese registro completo y no se deduce de la lista
\((e_i)\) por una identificación de nombres.


% BEGIN OWNER /Users/ruben/Documents/New project/output/EXTREMA_MEDIA_RAZON_RECIPROCIDAD_ES_EN_20260918/ARTICULO/ES/source/sections/registro_incidencias.tex
% Suplemento de trabajo. No sustituye registro_k.tex ni cierra por decreto E108.
% Procedencia y alcance: INFORME_RECUPERACION.md, misma carpeta.
\subsection{Evaluación incidencial anterior al registro dodecafásico}
\label{x_reciprocidad:subsec:registro-evaluacion-incidencial}

La recuperación de un registro a partir de sus diferencias cíclicas y su
carga total es una cuestión distinta de la evaluación de esas coordenadas
sobre el registro de incidencias. El primer problema se resuelve mediante el
extractor transversal de la proposición~\ref{x_reciprocidad:prop:k-transversal}.
Para el segundo, una presentación histórica del
registro proporciona tres recuentos enteros en cada ventana. Es útil
explicitar su composición y determinar exactamente la información que
requiere su evaluación.

Sea \(\mathscr L\) un libro de visitas y trazas orientadas previamente
generado. Cada visita conserva la ruta, la posición APP, la hoja, el
instante, la multiplicidad y la incidencia de frontera. La evaluación
aritmética produce el entero y su descomposición
\[
 n_t=r_t+9q_t,\qquad 1\leq r_t\leq9.
\]
El entero \(n_t\) se calcula a partir de la hoja de suma o producto; su
residuo no se recibe como un valor independiente. La distinción entre
incidencia de corona e incidencia interior permanece explícita para las
visitas de residuo \(9\).

En la ventana \(E_m\), la presentación por recuentos considera
\begin{align}
 A_m&=\sum_{v\in E_m}w(v)
             \mathbf1_{\{1,3,5,7\}}(r(v)),\\
 C_m&=\sum_{j\geq0}\bigl(\nu^-_{120,m}(j)-\nu^+_{120,m}(j)\bigr),\\
 V_m&=\sum_{v\in E_m}w(v)\mathbf1_{\{9\}}(r(v))\epsilon_\partial(v).
\end{align}
Aquí \(w(v)\) es la multiplicidad incidencial y
\(\epsilon_\partial(v)=1\) en la corona, \(-1\) en el interior.
Las trazas contadas por \(\nu^\pm\) deben enumerarse mediante la regla
correspondiente del artículo. La suma entera requiere soporte finito o una
construcción alternativa que determine su significado. El número de
incidencias cronológicas y la longitud reducida \(120(1+3j)\) tienen
tipos distintos; su relación exige la unidad de longitud del lector.

Se define \([x]_{10}\in\{0,\ldots,9\}\) como el representante
euclídeo y se conservan los cocientes
\[
 x=[x]_{10}+10\lfloor x/10\rfloor.
\]
El bloque incidencial es
\begin{equation}
 z_m=100[A_m]_{10}+10[C_m]_{10}+[V_m]_{10}.
 \label{x_reciprocidad:eq:registro-bloque-incidencial}
\end{equation}
Por composición se obtiene la aplicación
\begin{equation}
 \mathcal E_{\rm cnt}(\mathscr L)=
 \left((z_m-z_{m+3})_{m=1}^{12},
       (z_m-z_{m+4})_{m=1}^{12},\sum_{m=1}^{12}z_m\right).
\label{x_reciprocidad:eq:registro-composicion-incidencial}
\end{equation}
Los índices se leen cíclicamente módulo \(12\). En este apartado,
\((D_jz)_m=z_m-z_{m+j}\), para \(j=3,4\), y
\(Q(z)=\sum_{m=1}^{12}z_m\); por tanto, la composición anterior es
\((D_3z,D_4z,Q(z))\). Esta convención se conserva en
\S\ref{x_reciprocidad:subsec:k-transversal}, donde se demuestra la reconstrucción
integral completa.
Esta fórmula produce sus coordenadas a partir de los recuentos. No utiliza
\(K\), \(U\), \(\alpha\) ni las coordenadas transversales esperadas
como argumentos. Su identificación con
\(\mathcal E_{108}^{90,120}\) requiere componerla con la familia y
las incidencias efectivas del estado terminal; no se deduce únicamente
de la invertibilidad del sistema transversal.

\begin{proposition}[Información aritmética suficiente y necesaria]
\label{x_reciprocidad:prop:registro-36-residuos}
Sean \(N,N'\in\mathbb Z^{12\times3}\), ordenados por los recuentos
\((A,C,V)\). Para la aplicación definida por
\eqref{x_reciprocidad:eq:registro-bloque-incidencial}--\eqref{x_reciprocidad:eq:registro-composicion-incidencial},
\[
 \mathcal E_{\rm cnt}(N)=\mathcal E_{\rm cnt}(N')
 \quad\Longleftrightarrow\quad
 N-N'\in10\mathbb Z^{36}.
\]
En particular, los \(36\) residuos sectoriales determinan exactamente
la salida de \(25\) coordenadas. La conservación de los cocientes y
de la procedencia constituye información adicional del registro.
\end{proposition}

\begin{proof}
La igualdad módulo \(10\) produce los mismos bloques y, por tanto,
las mismas coordenadas transversales. Recíprocamente, si las coordenadas
transversales coinciden, \(D_3(z-z')=D_4(z-z')=0\). Los pasos
\(3\) y \(4\) generan el ciclo de longitud \(12\), de modo que
\(z-z'\) es constante. La igualdad de cargas totales anula esa
constante. Finalmente, la representación decimal de un entero entre
\(0\) y \(999\) tiene tres dígitos únicos. Se recuperan así los
tres residuos de cada ventana. La afirmación es una identidad algebraica
para todo \(N,N'\), no una inferencia obtenida de pruebas finitas.
\end{proof}

\subsection{Conjugación de los recuentos y términos de frontera}
\label{x_reciprocidad:subsec:registro-conjugacion-incidencial}

Sea \(\rho(m)=13-m\). Considérese una conjugación de incidencias
que invierte el orden de las ventanas, conserva la clase aritmética y la
clase de frontera de las visitas e intercambia los canales de las trazas.
Entonces
\[
 (A_m^\dagger,C_m^\dagger,V_m^\dagger)
 =(A_{\rho(m)},-C_{\rho(m)},V_{\rho(m)}).
\]
Si \(c_m=[C_m]_{10}\), la reducción decimal produce la identidad
\begin{equation}
 z_m^\dagger=z_{\rho(m)}
       +100\mathbf1_{\{c_{\rho(m)}\ne0\}}-20c_{\rho(m)}.
 \label{x_reciprocidad:eq:registro-conjugacion-decimal}
\end{equation}
En efecto, \([-C]_{10}=0\) cuando \([C]_{10}=0\), y
\([-C]_{10}=10-[C]_{10}\) en otro caso. La fórmula conserva
exactamente esa diferencia. La negación del canal no equivale a negar
el bloque decimal completo.

La aplicación de esta identidad a una involución concreta del TPK exige
verificar su acción sobre las incidencias. Para un lector de ocupaciones
evaluado antes del avance, la inversión de una ruta
\(\gamma=(x_0,\ldots,x_n)\) satisface
\[
 \sum_{k=0}^{n-1}f(x_{n-k})
 =\sum_{k=0}^{n-1}f(x_k)+f(x_n)-f(x_0).
\]
Los términos de los extremos se incorporan antes de reducir módulo
\(10\). En rutas cerradas se anulan; en ventanas abiertas pueden
alterar sus residuos. Esta corrección, ya presente en el desarrollo
bidireccional del transporte, especifica qué debe conservar el empalme
entre la ruta y los recuentos.

\begin{proposition}[Necesidad de la información no periódica del registro]
Supóngase que el observable de una familia fija de rutas satisface
\(o(t+54)=o(t)\), que la multiplicidad se conserva bajo ese retorno
y que el mismo lector local actúa en las ventanas de \(9\) instantes.
Entonces sus recuentos y bloques satisfacen
\[
 (A,C,V)_{m+6}=(A,C,V)_m,\qquad z_{m+6}=z_m.
\]
\end{proposition}

\begin{proof}
La traslación \(t\mapsto t+54\) lleva \(E_m\) biyectivamente
a \(E_{m+6}\), conserva cada incidencia contada y su multiplicidad.
La igualdad se transmite a la suma y a la reducción decimal.
\end{proof}

Por tanto, una presentación con \(12\) bloques distintos requiere que
el lector conserve alguna información que no factorice por ese observable
de período \(54\): selección de incidencias, orientación efectiva,
extremos, multiplicidad o memoria. La proposición delimita una reducción
inadmisible del estado; no afirma que la dinámica completa tenga
período \(54\) ni determina por sí misma cuál es su lector terminal.

% END OWNER


% BEGIN OWNER /Users/ruben/Documents/New project/output/EXTREMA_MEDIA_RAZON_RECIPROCIDAD_ES_EN_20260918/ARTICULO/ES/source/sections/censo_terminal_completo.tex
\subsection{Censo completo de cilindros y selección de la frontera terminal}
\label{x_reciprocidad:censo:terminal-completo}

La lectura terminal utiliza conjuntamente dos informaciones: la compatibilidad
de los cilindros regionales y un perfil de incidencia con canales y posiciones
etiquetados. En este apartado se construyen los catálogos completos, se
enumeran sus ternas y se demuestra el efecto de cada condición del lector.
El punto de partida son las salidas regionales ya producidas en la
sección~\ref{x_reciprocidad:sec:generacion}, no valores introducidos para definir el
transporte APP--TRIT--TPK. La reducción a una matriz pertenece a una lectura
posterior de esas salidas y conserva el orden de los tres canales.

El resultado es un teorema censal con estructura de partida explícita.
El perfil de Gram, pesos, distancias y márgenes se declara como dato del
lector terminal. La exhaustividad de la enumeración demuestra qué selecciona
ese perfil sobre los catálogos fijados; no demuestra, por sí sola, que el
perfil sea una ley universal ni que pueda reconstruirse desde la única terna
que finalmente lo satisface. Esta separación permite incorporar la prueba
finita completa sin invertir su genealogía.

\subsubsection{De las salidas regionales a los catálogos ternarios}

Sean \(x_\pi,x_e,x_\varphi\in[0,1)\) las partes fraccionarias de las
tres coordenadas regionales ya construidas. Se utilizan sus cilindros
publicados de mil posiciones decimales,
\[
 J_\chi^{(1000)}=
 \left[\frac{d_\chi}{10^{1000}},
       \frac{d_\chi+1}{10^{1000}}\right),
 \qquad \chi\in\{\pi,e,\varphi\}.
\]
Los enteros \(d_\chi\) son lecturas finitas de salida. La longitud mil
es suficiente para los enteros que siguen, como se comprueba mediante las
dos cotas de cada intervalo; no es una profundidad supuesta infinita.
Definimos
\begin{equation}
 p=30,\quad r=1050,\quad L=p+r=1080,\quad k=171,
 \qquad D=3^L,\quad Q=1000^k,
 \label{x_reciprocidad:censo:parametros}
\end{equation}
y
\[
 P_\chi=\lfloor3^{30}x_\chi\rfloor,\qquad
 T_\chi=\lfloor Qx_\chi\rfloor.
\]
Aquí \(k\) cuenta tríadas decimales y no designa el registro dodecafásico
\(K\). Se tiene \(1000^{171}\leq3^{1080}<1000^{172}\).

\begin{lemma}[Extracción estable de un entero]
\label{x_reciprocidad:censo:estabilidad}
Para \(d,m,h\in\mathbb Z\), con \(m,h>0\), el valor de
\(\lfloor mx\rfloor\) es constante en \([d/h,(d+1)/h)\) si y sólo si
\begin{equation}
 \left\lfloor\frac{md}{h}\right\rfloor
 =\left\lfloor\frac{m(d+1)-1}{h}\right\rfloor.
 \label{x_reciprocidad:censo:prueba-estabilidad}
\end{equation}
Cuando se cumple, cualquiera de los dos miembros da ese valor.
\end{lemma}
\begin{proof}
El mínimo de \(mx\) se alcanza en el extremo izquierdo. Su supremo
es \(m(d+1)/h\) y no se alcanza. El mayor entero que puede tomar
\(\lfloor mx\rfloor\) es
\(\lceil m(d+1)/h\rceil-1\), igual al segundo miembro porque
el numerador y el denominador son enteros. La igualdad de los extremos
del intervalo de valores prueba la afirmación.
\end{proof}

La prueba de estabilidad se cumple en los tres canales para
\(m=3^{30},Q,3^{1080}\). La última elección se utilizará únicamente
para una comparación posterior, no para seleccionar el perfil. Las
primeras treinta posiciones recuperadas son precisamente las cinco
etapas de seis trits del transporte regional:
\begin{center}\small
\begin{tabular}{@{}c l r@{}}\toprule
Canal&Prefijo de treinta trits&\(P_\chi\)\\\midrule
\(\pi\)&\texttt{010211012222010211002111110221}&29152671743887\\
\(e\)&\texttt{201101121221102011012222102011}&147887858824447\\
\(\varphi\)&\texttt{121200112202121020010210010200}&127247717616687\\\bottomrule
\end{tabular}
\end{center}
La coincidencia conserva la prolongación ya construida; la lectura decimal
no se usa para sustituir retroactivamente sus matrices de transporte.

Fijado un canal, una cola de \(r\) trits se representa por
\(0\leq R<3^r\). Su cilindro y el cilindro de las primeras \(k\)
tríadas son
\begin{equation}
 I_{\chi,R}=\left[\frac{P_\chi3^r+R}{D},
                       \frac{P_\chi3^r+R+1}{D}\right),\qquad
 J_{\chi,T}=\left[\frac{T_\chi}{Q},\frac{T_\chi+1}{Q}\right).
 \label{x_reciprocidad:censo:dos-cilindros}
\end{equation}
El catálogo incluye todas las colas para las que estos dos cilindros
tienen intersección no vacía.

\begin{proposition}[Catálogo por intersección exacta]
\label{x_reciprocidad:censo:catalogos}
Las colas compatibles son los enteros del intervalo \([a_\chi,b_\chi]\),
con
\begin{align}
 a_\chi&=\max\left(0,
       \left\lfloor\frac{T_\chi D}{Q}\right\rfloor-P_\chi3^r\right),
       \label{x_reciprocidad:censo:extremo-a}\\
 b_\chi&=\min\left(3^r-1,
       \left\lfloor\frac{(T_\chi+1)D-1}{Q}\right\rfloor-P_\chi3^r\right).
       \label{x_reciprocidad:censo:extremo-b}
\end{align}
Para los cilindros regionales fijados, los recortes no actúan y se obtiene
\begin{equation}
\begin{array}{c|r|r|r}
 \chi&a_\chi\bmod729&b_\chi\bmod729&b_\chi-a_\chi+1\\\hline
 \pi&67&262&196\\
 e&584&51&197\\
 \varphi&117&312&196
\end{array}
\label{x_reciprocidad:censo:tabla-catalogos}
\end{equation}
\end{proposition}
\begin{proof}
Escribamos \(N=P_\chi3^r+R\). Dos intervalos semiabiertos de
\eqref{x_reciprocidad:censo:dos-cilindros} se cortan exactamente cuando
\[
 NQ<(T_\chi+1)D,\qquad (N+1)Q>T_\chi D.
\]
La segunda desigualdad equivale a
\(N\geq\lfloor T_\chi D/Q\rfloor\); la primera, a
\(N\leq\lfloor((T_\chi+1)D-1)/Q\rfloor\). Restar
\(P_\chi3^r\) y conservar \(0\leq R<3^r\) prueba las fórmulas.
La evaluación por divisiones enteras de los \(T_\chi\) extraídos por
\eqref{x_reciprocidad:censo:prueba-estabilidad} da la tabla. Como \(r\geq6\),
\(P_\chi3^r\equiv0\pmod{729}\), de modo que el residuo del primer
extremo también se obtiene directamente de
\(\lfloor T_\chi D/Q\rfloor\bmod729\).
\end{proof}

Los últimos seis trits de una cola forman la palabra
\(\nu_6(R\bmod729)\), donde \(\nu_6\) es la escritura ternaria
de longitud seis, incluidos sus ceros iniciales. Por la tabla, cada
catálogo tiene menos de \(729\) elementos consecutivos, de manera que
esta reducción es inyectiva sobre él. Así, los tres catálogos de frontera
quedan descritos completamente, sin elegir las ternas finales:
\begin{equation}
 \begin{split}
 \mathcal B_\pi&=\{\nu_6(67+i):0\leq i<196\},\\
 \mathcal B_e&=\{\nu_6((584+j)\bmod729):0\leq j<197\},\\
 \mathcal B_\varphi&=\{\nu_6(117+\ell):0\leq\ell<196\}.
 \end{split}
 \label{x_reciprocidad:censo:fronteras-explicitas}
\end{equation}
Exigir que el extremo izquierdo de \(I_{\chi,R}\) pertenezca a
\(J_{\chi,T}\) sería otra condición: perdería el primer cilindro
que corta la frontera y daría \(195,196,195\). El censo utiliza la
intersección de los cilindros completos, de acuerdo con su definición.

\subsubsection{Perfil declarado e indicadores de compatibilidad}

Para \(u,v\in\mathbb F_3^6\), sean
\(g(u,v)=\sum u_iv_i\in\mathbb F_3\),
\(n(u)=g(u,u)\), \(w(u)=\#\{i:u_i\ne0\}\) y
\(h(u,v)=\#\{i:u_i\ne v_i\}\). Las dos últimas funciones toman
valores enteros; las dos primeras, valores en \(\mathbb F_3\).
El perfil local del lector, en orden \((\pi,e,\varphi)\), es
\begin{equation}
 \begin{split}
 (g_{\pi e},g_{\pi\varphi},g_{e\varphi})&=(2,2,0),\qquad
 (n_\pi,n_e,n_\varphi)=(0,0,0),\\
 (w_\pi,w_e,w_\varphi)&=(3,3,3),\qquad
 (h_{\pi e},h_{\pi\varphi},h_{e\varphi})=(2,5,3).
 \end{split}
 \label{x_reciprocidad:censo:perfil-local}
\end{equation}
Las palabras se ordenan como en \eqref{x_reciprocidad:censo:fronteras-explicitas}.
Usaremos \([E]\in\{0,1\}\) para el indicador de una condición \(E\).
Para cada par de canales \(\chi,\psi\), la matriz rectangular
\(A_{\chi\psi}\) tiene entrada \([g(u,v)=g_{\chi\psi}]\).
La matriz \(A^{h}_{\chi\psi}\) multiplica esa entrada por
\([h(u,v)=h_{\chi\psi}]\). Sean \(D_\chi^n\) la matriz diagonal
de \([n(u)=0]\), y \(D_\chi^{nw}\) la de
\([n(u)=0][w(u)=3]\).

\begin{theorem}[Enumeración exhaustiva por sumas finitas]
\label{x_reciprocidad:censo:teorema-conteo}
El producto de los tres catálogos contiene \(7\,567\,952\) ternas.
Las condiciones sucesivas de \eqref{x_reciprocidad:censo:perfil-local} dejan
\begin{equation}
 7\,567\,952\longrightarrow275\,794\longrightarrow6235
 \longrightarrow3127\longrightarrow331.
 \label{x_reciprocidad:censo:cadena-local}
\end{equation}
Estos cuatro conteos son, respectivamente,
\begin{align}
 N_g&=\operatorname{tr}(A_{\pi e}A_{e\varphi}A_{\varphi\pi}),
 \nonumber\\
 N_{gn}&=\operatorname{tr}(D_\pi^n A_{\pi e}
             D_e^n A_{e\varphi}D_\varphi^n A_{\varphi\pi}),
 \nonumber\\
 N_{gnw}&=\operatorname{tr}(D_\pi^{nw} A_{\pi e}
             D_e^{nw} A_{e\varphi}D_\varphi^{nw} A_{\varphi\pi}),
 \nonumber\\
 N_{gnwh}&=\operatorname{tr}(D_\pi^{nw} A^h_{\pi e}
             D_e^{nw} A^h_{e\varphi}D_\varphi^{nw} A^h_{\varphi\pi}).
 \label{x_reciprocidad:censo:trazas}
\end{align}
Todos los productos y trazas de esta fórmula se calculan en los enteros,
no en el cuerpo de tres elementos.
\end{theorem}
\begin{proof}
El tamaño del producto es \(196\cdot197\cdot196\). Al expandir
una traza de \eqref{x_reciprocidad:censo:trazas}, sus índices recorren exactamente una
terna de palabras. El sumando vale uno si satisface todas las condiciones
representadas y cero en otro caso. Cada terna se cuenta una vez.
Las matrices diagonales añaden las condiciones individuales; los factores
rectangulares añaden las condiciones por pares. Esto demuestra que las
trazas son precisamente los cardinales anunciados.

Para explicitar su evaluación, fijados los índices \(i,j\) de los dos
primeros canales, se forman los conjuntos de índices \(\ell\) del tercero
que cumplen cada condición por pares. Se intersectan sus indicadores,
se añaden los indicadores individuales pertinentes y se suma su cardinal.
La suma final recorre \(0\leq i<196\), \(0\leq j<197\).
La sustitución de las palabras de \eqref{x_reciprocidad:censo:fronteras-explicitas} y
del perfil \eqref{x_reciprocidad:censo:perfil-local} da, en ese orden,
\(275794,6235,3127,331\). El suplemento registra las cuatro contribuciones
de cada uno de los \(196\) índices \(i\) y las \(331\) ternas completas;
su programa evalúa estas sumas con indicadores binarios exactos y contrasta
escalarmente cada superviviente. No interviene una muestra ni una tolerancia.
\end{proof}

\subsubsection{Incidencia etiquetada y carga de orientación}

Para una terna superviviente, sea \(B\in\mathbb F_3^{3\times6}\) su
matriz de filas ordenadas. El segundo perfil exige
\begin{equation}
 \operatorname{rank}_{\mathbb F_3}B=3,\qquad
 w_{\mathrm{col}}(B)=(1,1,2,2,3,0),\qquad
 q_{\partial}(B)=(1,2,2,1,2,0),
 \label{x_reciprocidad:censo:perfil-columnas}
\end{equation}
donde \(w_{\mathrm{col}}\) cuenta los elementos no nulos de cada
columna y \(q_{\partial}\) suma sus entradas módulo tres. Es un perfil
de columnas etiquetadas; permutarlas sin transportar el perfil define otro
lector. Denotemos por \(\mathcal F\) el conjunto de las \(331\) ternas
del teorema anterior.

\begin{proposition}[Reducción del censo al par terminal]
\label{x_reciprocidad:censo:dos-matrices}
Las sumas de indicadores sobre \(\mathcal F\), al exigir sucesivamente
rango, pesos de columna y sumas de columna, son
\begin{equation}
 \sum_{B\in\mathcal F}[\operatorname{rank}B=3]=331,
 \qquad N_{\mathrm{rango,pesos}}=18,
 \qquad N_{\mathrm{rango,pesos,sumas}}=2.
 \label{x_reciprocidad:censo:conteo-columnas}
\end{equation}
Con índices de catálogo comenzando en uno, las dos ternas finales son
\((120,197,157)\) y \((129,197,148)\), y sus matrices son
\[
 B_0=\begin{pmatrix}0&2&0&2&2&0\\0&0&1&2&2&0\\1&0&1&0&1&0\end{pmatrix},
 \qquad
 B_1=\begin{pmatrix}0&2&1&0&2&0\\0&0&1&2&2&0\\1&0&0&2&1&0\end{pmatrix}.
\]
\end{proposition}
\begin{proof}
En cada terna del conjunto ya enumerado se calcula el rango por eliminación
en \(\mathbb F_3\), se cuentan los elementos no nulos de las seis
columnas y se suman éstas módulo tres. Los productos de los indicadores
de \eqref{x_reciprocidad:censo:perfil-columnas} dan \eqref{x_reciprocidad:censo:conteo-columnas}.
La sustitución de los dos índices finales en
\eqref{x_reciprocidad:censo:fronteras-explicitas} da las matrices expuestas. La lista
completa de supervivientes y los filtros exactos del suplemento permiten
reproducir la eliminación sin presuponer esas dos matrices como dominio.
\end{proof}

La orientación \(\sigma=(1,1,-1)=(1,1,2)\in\mathbb F_3^3\)
selecciona la recta
\(\langle\sigma\rangle=\{(0,0,0),(1,1,2),(2,2,1)\}\).
Las cargas de filas de \(B_0,B_1\) son \((0,2,0)\) y \((2,2,1)\);
sólo la segunda pertenece a esa recta. Se recupera así el último paso
del selector de la sección~\ref{x_reciprocidad:subsec:k-terminal}, ahora precedido por
la construcción y enumeración de su dominio completo. La lectura estable
de \(\lfloor3^{1080}x_\chi\rfloor\), efectuada después de los filtros,
da las fronteras \texttt{021020}, \texttt{001220}, \texttt{100210},
con los mismos rangos \((129,197,148)\). Este cotejo posterior no define
el perfil ni participa en las sumas de indicadores.

\subsubsection{Información conservada y alcance del resultado}

El procedimiento conserva el prefijo de treinta trits, la cola compatible,
el orden regional, la posición de cada trit y la orientación. La reducción
módulo \(729\) conserva aquí la identidad de cada candidato porque el
intervalo de colas tiene longitud menor que \(729\); esa inyectividad
no se afirma para una historia arbitraria. El censo de \(331\) demuestra
también que el perfil local aislado no selecciona una sola matriz: los
márgenes etiquetados y la carga aportan condiciones efectivas adicionales.

Se ha recuperado, por tanto, una enumeración finita íntegra y su selector
sobre cilindros de salida, con perfil de lectura explícito. La construcción
del registro firmado utiliza además su lector de memoria y sus canales,
como se especifica a continuación; la matriz visible no se identifica con
ese registro por igualdad de número de componentes. Tampoco este conteo
demuestra la universalidad del perfil, genera de nuevo las coordenadas
regionales ni decide una afirmación sobre alfa o sobre el continuo completo.

% END OWNER


\subsection{Dos lecturas del estado terminal}
\label{x_reciprocidad:subsec:k-terminal}

Denotemos por \(x_{\mathrm{term}}\) el estado producido por la composición
de selección, transporte y actualización hasta el corte terminal. La
escritura operatoria es
\[
 x_{\mathrm{term}}=
 (\mathcal U_{108}\circ\cdots\circ\mathcal U_1)(x_{\mathrm{seed}}),
 \qquad \mathcal U_t=\mathsf{Upd}_t\circ\mathsf{Tra}_t\circ\mathsf{Sel}_t.
\]
Sus dos lecturas pertinentes son
\[
 x_{\mathrm{term}}
 \xrightarrow{(\pi_{\mathrm{term}},R_{\mathrm{sgn}})}
 (B_{\mathrm{term}},U).
\]
La matriz \(B_{\mathrm{term}}\) retiene la lectura ternaria visible; el
vector \(U\) conserva otra coordenada de la misma historia. No se inserta
entre ellos una flecha \(B_{\mathrm{term}}\to U\).

El selector terminal opera sobre los tres catálogos regionales
de tamaños 196, 197 y 196, construidos en la
sección~\ref{x_reciprocidad:censo:terminal-completo}. Las firmas locales reducen sus
\(196\cdot197\cdot196=7\,567\,952\) ternas a 331, y el margen de
columnas \(q_{\partial}=(1,2,2,1,2,0)\) deja dos matrices:
\[
 B_0=\begin{pmatrix}0&2&0&2&2&0\\0&0&1&2&2&0\\1&0&1&0&1&0\end{pmatrix},
 \qquad
 B_1=\begin{pmatrix}0&2&1&0&2&0\\0&0&1&2&2&0\\1&0&0&2&1&0\end{pmatrix}.
\]
La orientación de filas es \(\sigma=(1,1,-1)=(1,1,2)\) en
\(\mathbb F_3^3\). Las cargas de filas de las candidatas son
\((0,2,0)\) y \((2,2,1)\). Sólo la segunda pertenece a
\(\langle\sigma\rangle=\{(0,0,0),(1,1,2),(2,2,1)\}\); por tanto
\(B_{\mathrm{term}}=B_1\). Esta comprobación prueba el último paso del
selector a partir del censo declarado. No convierte el margen de columnas
en una consecuencia de la carga de filas, ni reemplaza la enumeración
anterior por la inspección de esas dos matrices.

El lector firmado tiene la composición tipada
\begin{equation}
 R_{\mathrm{sgn}}
 =\Pi_H\circ\mathcal E_{108}^{90,120}\circ\mathcal R_{12}:
 \mathcal X_{\mathrm{TPK}}^{\mathrm{term,mem}}
 \longrightarrow\mathbb Z^{12}.
 \label{x_reciprocidad:eq:k-lector-firmado}
\end{equation}
El registro de canales determina
\[
 \mathcal E_{108}^{90,120}(\mathcal R_{12}(x_{\mathrm{term}}))
 =(b^{90},b^{120},Q_{\mathrm{TPK}}),
\]
con las coordenadas
\begin{align*}
 b^{90}={}&(-495,-116,-684,108,601,-90,\\
            &\hspace{19mm}-173,-88,313,560,-397,461),\\
 b^{120}={}&(-425,-281,-481,671,-255,30,\\
             &\hspace{19mm}475,-543,680,251,6,-128),\\
 Q_{\mathrm{TPK}}={}&6263.
\end{align*}
Estas cantidades se utilizan como coordenadas del registro previo, no
como parámetros que se ajusten mediante alfa. La reconstrucción demostrada
en los apartados siguientes comienza exactamente en esta salida tipada.
La caracterización integral del extractor, su inversa y la composición
de contribuciones se desarrollan en la sección~\ref{x_reciprocidad:img09:seccion},
en particular en el teorema~\ref{x_reciprocidad:img09:imagen} y la
proposición~\ref{x_reciprocidad:img09:acumulacion}. La determinación regional
del descriptor y la selección terminal se exponen en el
capítulo~\ref{chap:despliegue-selector-k}. Cada construcción conserva su
dominio: el selector actúa sobre los repertorios generados y la inversión
actúa sobre los canales firmados declarados aquí. El rastro documental del
extractor queda registrado en la procedencia técnica del paquete.

\subsection{Construcción del registro firmado mediante los canales orientados}
\label{x_reciprocidad:subsec:k-pi-h}

El lector armónico \(\Pi_H\) puede escribirse sin utilizar ni \(K\) ni
alfa. Para evitar confundir sus sumas orbitales con los bloques de alfa,
las denotaremos por \(s_1,s_2,s_3\). Sean
\[
 B_r=\sum_{j=0}^{3}b^{120}_{r+3j},\qquad
 d_{r,j}=b^{90}_{r+3j},\qquad1\leq r\leq3,
\]
y
\begin{equation}
 s_1=\frac{Q_{\mathrm{TPK}}+2B_1+B_2}{3},\qquad
 s_2=s_1-B_1,\qquad s_3=s_2-B_2.
 \label{x_reciprocidad:eq:k-sumas-armonicas}
\end{equation}
La división por tres pertenece al lector de las tres órbitas. Su
integralidad se comprueba sobre las salidas compatibles del registro; no
se supone para un elemento arbitrario de \(\mathbb Z^{25}\).

Los tres bloques firmados son
\begin{equation}
 U^{(r)}=(s_r,d_{r,1}-d_{r,3},d_{r,0}+d_{r,2},d_{r,0}-d_{r,2}).
 \label{x_reciprocidad:eq:k-bloques-firmados}
\end{equation}
La sustitución entera da
\[
 (B_1,B_2,B_3)=(972,-1073,101),\qquad
 (s_1,s_2,s_3)=(2378,1406,2479),
\]
y, por tanto,
\begin{align*}
 U^{(1)}&=(2378,-452,-668,-322),\\
 U^{(2)}&=(1406,998,-204,-28),\\
 U^{(3)}&=(2479,-551,-371,-997).
\end{align*}
Intercalando por columnas se obtiene
\begin{equation}
 U=(2378,1406,2479,-452,998,-551,
       -668,-204,-371,-322,-28,-997).
 \label{x_reciprocidad:eq:k-u-firmado}
\end{equation}
La presencia de valores negativos y de componentes superiores a 999
muestra que \(U\) no es una palabra de dígitos en base mil. Tampoco es
\(U_6\Vert U_6\), concatenación catalogal de período seis. La diferencia
es de dominio y de información conservada, no solamente de valores.

El orden constructivo hasta aquí es
\[
 x_{\mathrm{term}}\longrightarrow\mathcal R_{12}
 \longrightarrow(b^{90},b^{120},Q_{\mathrm{TPK}})
 \longrightarrow U.
\]
La inversión que sigue produce el registro dodecafásico. Que después podamos recuperar
los canales desde el registro dodecafásico constituye un control de reversibilidad; no
invierte este orden causal.

\subsection{Hadamard por órbitas y reconstrucción entera}
\label{x_reciprocidad:subsec:k-hadamard}

En el ciclo de doce posiciones, el paso tres determina tres órbitas:
\[
 (1,4,7,10),\qquad(2,5,8,11),\qquad(3,6,9,12).
\]
La separación de cuatro sectores de cada órbita es la que, una vez fijados
origen y orientación, admite la lectura angular de noventa grados. La
transformación no actúa sobre tres cuaternas contiguas del orden natural.

Sea
\[
 H_4=\begin{pmatrix}
 1&1&1&1\\1&1&-1&-1\\1&-1&1&-1\\1&-1&-1&1
 \end{pmatrix}.
\]
Si \(P_{\mathrm{orb}}\) reordena por las tres órbitas indicadas,
definimos
\[
 \mathcal H_{12}=P_{\mathrm{orb}}^{-1}
 (H_4\oplus H_4\oplus H_4)P_{\mathrm{orb}}.
\]

\begin{lemma}[Inversión de Hadamard e integralidad]
\label{x_reciprocidad:lem:k-hadamard}
Se tiene
\[
 H_4^2=4I_4,\qquad H_4^{-1}=\frac14H_4,\qquad\det H_4=-16.
\]
Para \(u\in\mathbb Z^4\), la preimagen \(H_4^{-1}u\) es entera si y
sólo si \(H_4u\equiv0\pmod4\).
\end{lemma}

\begin{proof}
Las filas son ortogonales y tienen norma cuadrada cuatro; la matriz es
simétrica. Esto da \(H_4^2=4I_4\) y la inversa. Una eliminación entera,
o la expansión del determinante, da el signo negativo y el valor \(-16\).
La condición de integralidad es exactamente la divisibilidad por cuatro
de las cuatro coordenadas de \(H_4u\).
\end{proof}

\begin{definition}[Transformación integral dodecafásica]
La subred de registros firmados y su coordenatización integral son
\[
 \mathcal L_H=\{u\in\ZZ^{12}:\mathcal H_{12}u\equiv0\pmod4\},
 \qquad \mathscr K(u)=\tfrac14\mathcal H_{12}u.
\]
La aplicación \(\mathscr K:\mathcal L_H\to\ZZ^{12}\) es un isomorfismo
de grupos abelianos, con inversa \(k\mapsto\mathcal H_{12}k\).
El registro canónico \(K\) es la imagen del registro firmado generado
por las operaciones anteriores, con origen de fase y orientación fijados.
\end{definition}

\begin{proof}
La congruencia define exactamente la integralidad de \(\mathscr K(u)\).
La identidad \(\mathcal H_{12}^2=4I\) demuestra las dos composiciones
inversas. En particular, todo \(k\in\ZZ^{12}\) tiene preimagen
\(\mathcal H_{12}k\in\mathcal L_H\).
\end{proof}

\begin{proposition}[Registro integral canónico]
\label{x_reciprocidad:prop:k-sello}
La inversión \(K^{(r)}=H_4U^{(r)}/4\) de los bloques de
\eqref{x_reciprocidad:eq:k-bloques-firmados} produce
\begin{align*}
 K^{(1)}&=(234,729,621,794),\\
 K^{(2)}&=(543,659,58,146),\\
 K^{(3)}&=(140,824,914,601).
\end{align*}
En el orden natural, el registro dodecafásico es
\begin{equation}
 K=(234,543,140,729,659,824,621,058,914,794,146,601).
 \label{x_reciprocidad:eq:k-vector}
\end{equation}
Es la única preimagen racional de \(U\), y pertenece a
\(\{0,\ldots,999\}^{12}\).
\end{proposition}

\begin{proof}
Los tres productos antes de dividir son
\begin{align*}
 H_4U^{(1)}&=(936,2916,2484,3176),\\
 H_4U^{(2)}&=(2172,2636,232,584),\\
 H_4U^{(3)}&=(560,3296,3656,2404).
\end{align*}
Todas las coordenadas son divisibles por cuatro. Dividir y deshacer la
permutación de órbitas da \eqref{x_reciprocidad:eq:k-vector}. La unicidad procede de la
invertibilidad racional, no de una búsqueda entre candidatos próximos a
un valor físico.
\end{proof}

La integralidad también puede formularse en términos de redes. Los
divisores determinantes de \(H_4\) son \(1,2,4,16\): las entradas
tienen máximo común divisor uno; los menores de orden dos son pares y
alguno tiene módulo dos; los de orden tres son múltiplos de cuatro y
alguno tiene módulo cuatro. De ahí se deduce
\[
 \operatorname{SNF}(H_4)=\operatorname{diag}(1,2,2,4),
\]
y para la transformación global,
\[
 \operatorname{coker}\mathcal H_{12}
 \cong(\mathbb Z/2\mathbb Z)^6\oplus(\mathbb Z/4\mathbb Z)^3.
\]
El índice de la imagen es \(16^3=4096\). Así, no todo vector firmado
entero admite una preimagen entera; el vector obtenido satisface las
congruencias precisas que la inversión exige. El factor \(1/4\) está
forzado por la matriz, y no tiene el estatuto de un coeficiente libre.

\subsection{Diferencias cíclicas y componente uniforme}
\label{x_reciprocidad:subsec:k-transversal}

Con índices cíclicos módulo doce, sean
\[
 (D_3z)_m=z_m-z_{m+3},\qquad(D_4z)_m=z_m-z_{m+4},
 \qquad Q(z)=\sum_{m=1}^{12}z_m.
\]
El primer operador compara sectores separados por tres posiciones; el
segundo, por cuatro. Los nombres de canal 90 y 120 se refieren aquí a
esas separaciones en el ciclo orientado, no a una identificación con un
operador físico de otro dominio.

\begin{proposition}[Completitud del sistema transversal]
\label{x_reciprocidad:prop:k-transversal}
El operador
\[
 \mathcal D=(D_3,D_4,Q):\mathbb Z^{12}\longrightarrow\mathbb Z^{25}
\]
tiene rango doce y factores invariantes no nulos
\[
 \operatorname{diag}(1,\ldots,1,12),
\]
con once unidades. En particular, sus salidas determinan un único registro dodecafásico.
\end{proposition}

\begin{proof}
Si \(D_3z=D_4z=0\), el vector es constante a lo largo de ambos pasos.
Como \(3\) y \(4\) generan el ciclo módulo doce,
\(\ker D_3\cap\ker D_4=\langle\mathbf1\rangle\). La carga total
elimina esta última libertad porque \(Q(\mathbf1)=12\).

Para la afirmación integral, las filas de \((D_3,D_4)\) son incidencias
orientadas de un grafo conexo. Un árbol generador proporciona once
factores invariantes unitarios. Todo menor de orden doce debe contener la
fila \(Q\). Sumando las primeras once columnas a la última, esta última
se anula en las filas de incidencia y vale doce en la fila \(Q\). Todos
los menores máximos son divisibles por doce y el árbol produce uno de
módulo exactamente doce. El último factor es, por tanto, doce.
\end{proof}

La comprobación sobre \eqref{x_reciprocidad:eq:k-vector} da
\[
 D_3K=b^{90},\qquad D_4K=b^{120},\qquad Q(K)=6263.
\]
Las fórmulas \eqref{x_reciprocidad:eq:k-sumas-armonicas} y
\eqref{x_reciprocidad:eq:k-bloques-firmados} reconstruyen a su vez \(U\) desde esas
diferencias. Se cierra así un triángulo de representaciones exactas del
mismo objeto: registro transversal, vector firmado y registro dodecafásico. La igualdad
de resultados no convierte sus codominios en uno solo; demuestra que
las aplicaciones declaradas recuperan la información pertinente.


% BEGIN OWNER /Users/ruben/Documents/New project/output/EXTREMA_MEDIA_RAZON_RECIPROCIDAD_ES_EN_20260918/ARTICULO/ES/source/sections/registro_imagen_integral.tex
% Formalización nueva integrada en la revisión III; originales preservados.
% Propietarios: U016, ecuaciones extractor-diferencias-k y lector-armonico-*;
% registro_k.tex, subsecciones k-pi-h, k-hadamard y k-transversal.
\subsection{Compatibilidad integral y determinación del registro transversal}
\label{x_reciprocidad:img09:seccion}

El dominio de esta construcción es la publicación del registro de una
historia APP--TRIT--TPK. Se distinguen la evaluación de los eventos del
registro, la compatibilidad de sus canales y la inversión de sus
coordenadas. El resultado siguiente da un criterio integral completo para
el segundo nivel y una factorización explícita para conectar el primero con él.
La demostración utiliza los operadores del registro y establece sus
relaciones para todo elemento de los dominios enteros indicados.

\subsubsection{Caracterización por incrementos adyacentes}

Los índices pertenecen a $\mathbb Z/12\mathbb Z$. Sean
\[
 (Sz)_i=z_{i+1},\qquad D_3=I-S^3,\qquad D_4=I-S^4,
 \qquad Q(z)=\sum_{i=0}^{11}z_i.
\]
Se conserva el extractor transversal
\[
 \mathcal D:\mathbb Z^{12}\longrightarrow\mathbb Z^{25},
 \qquad k\longmapsto(D_3k,D_4k,Q(k)).
\]
Para una terna de coordenadas enteras $(b,c,q)$, escribimos
\begin{equation}
 w_i=c_i-b_{i+1},\qquad
 P_0=0,\quad P_i=\sum_{j=0}^{i-1}w_j\ (1\leq i\leq12),
 \qquad T(w)=\sum_{i=0}^{11}P_i.
 \label{x_reciprocidad:img09:incremento}
\end{equation}

\begin{theorem}[Imagen integral del extractor transversal]
\label{x_reciprocidad:img09:imagen}
Una terna $(b,c,q)\in\mathbb Z^{25}$ pertenece a
$\operatorname{im}\mathcal D$ si y sólo si satisface
\begin{align}
 b_i&=w_i+w_{i+1}+w_{i+2}&& (0\leq i<12),
 \label{x_reciprocidad:img09:relaciones}\\
 \sum_{i=0}^{11}w_i&=0,
 \label{x_reciprocidad:img09:cierre}\\
 q+T(w)&\equiv0\pmod {12}.
 \label{x_reciprocidad:img09:congruencia}
\end{align}
Su única preimagen entera es
\begin{equation}
 k_i=\frac{q+T(w)}{12}-P_i,\qquad0\leq i<12.
 \label{x_reciprocidad:img09:inversa}
\end{equation}
Las doce ecuaciones de \eqref{x_reciprocidad:img09:relaciones} y la ecuación
\eqref{x_reciprocidad:img09:cierre} son linealmente independientes sobre $\mathbb Q$.
\end{theorem}
\begin{proof}
Para una salida de $\mathcal D$ se tiene
\[
 c_i-b_{i+1}=(k_i-k_{i+4})-(k_{i+1}-k_{i+4})=k_i-k_{i+1}.
\]
La suma de tres incrementos da $b_i$; la suma cíclica de todos ellos
es cero. Además, $P_i=k_0-k_i$, por lo que
$q=12k_0-T(w)$. Se obtienen todas las condiciones y la fórmula inversa.

Recíprocamente, \eqref{x_reciprocidad:img09:congruencia} hace entera la fórmula
\eqref{x_reciprocidad:img09:inversa}; \eqref{x_reciprocidad:img09:cierre} permite prolongarla
cíclicamente. Sus diferencias adyacentes son $w_i$. Por tanto,
$D_3k=b$ y
\[
 (D_4k)_i=w_i+w_{i+1}+w_{i+2}+w_{i+3}
 =w_i+b_{i+1}=c_i.
\]
La suma de la inversa da $Q(k)=q$. La misma fórmula prueba la unicidad.
Para la independencia, el cambio entero invertible
$(b,c)\mapsto(b,w=c-Sb)$ convierte las doce primeras relaciones en
$b=(I+S+S^2)w$. Cada una despeja una coordenada diferente de $b$.
La condición $\sum w_i=0$ es una relación adicional independiente.
\end{proof}

La imagen racional tiene dimensión doce. Su estructura integral añade
una congruencia de orden doce: ésta es la manifestación constructiva del
último factor de Smith de $\mathcal D$. Los once incrementos
$w_0,\ldots,w_{10}$ y la carga $q$ son coordenadas suficientes, con
$w_{11}=-\sum_{i=0}^{10}w_i$ y la congruencia
\eqref{x_reciprocidad:img09:congruencia}. Las veinticinco coordenadas originales
conservan esos mismos doce grados de libertad mediante trece relaciones.

\subsubsection{Conmutatividad con la lectura armónica}

Sea $\mathcal H_{12}$ la transformación de Hadamard del registro,
con órbitas $(r,r+3,r+6,r+9)$ para $r=0,1,2$. Su bloque es
\[
 H_4=\begin{pmatrix}
 1&1&1&1\\1&1&-1&-1\\1&-1&1&-1\\1&-1&-1&1
 \end{pmatrix},\qquad H_4^2=4I_4.
\]
Para $(b,c,q)$ compatible, se definen
\[
 B_r=\sum_{j=0}^3c_{r+3j},\qquad
 A_0=\frac{q+2B_0+B_1}{3},\quad A_1=A_0-B_0,
 \quad A_2=A_1-B_1.
\]
El lector $\Pi_H$ de \eqref{x_reciprocidad:eq:k-bloques-firmados} tiene los bloques
\[
 (\Pi_H(b,c,q))^{(r)}=
 (A_r,b_{r+3}-b_{r+9},b_r+b_{r+6},b_r-b_{r+6}).
\]

\begin{proposition}[Concordancia de las dos reconstrucciones]
\label{x_reciprocidad:img09:triangulo}
Sobre $\operatorname{im}\mathcal D$ se cumple
\[
 \Pi_H\mathcal D=\mathcal H_{12},\qquad
 \mathcal D^{-1}=\tfrac14\mathcal H_{12}\Pi_H.
\]
La imagen de $\Pi_H$ restringida a ese dominio es exactamente
\[
 \mathcal L_H=
 \{u\in\mathbb Z^{12}:\mathcal H_{12}u\equiv0\pmod4\}.
\]
\end{proposition}
\begin{proof}
Para $k=\mathcal D^{-1}(b,c,q)$, sea
$a_r=\sum_{j=0}^3k_{r+3j}$. El desplazamiento de cuatro posiciones
lleva cada órbita a la siguiente, y por ello $B_r=a_r-a_{r+1}$.
Junto con $q=a_0+a_1+a_2$, estas identidades dan $A_r=a_r$.
En una órbita, escribamos $(x_0,x_1,x_2,x_3)$ para las coordenadas
de $k$ y $d_j=x_j-x_{j+1}$. Entonces
\[
 d_1-d_3=x_0+x_1-x_2-x_3,\quad
 d_0+d_2=x_0-x_1+x_2-x_3,\quad
 d_0-d_2=x_0-x_1-x_2+x_3.
\]
Son las tres filas restantes de $H_4k^{(r)}$. La inversión y la
caracterización integral se siguen de $\mathcal H_{12}^2=4I$.
\end{proof}

La condición $\Pi_H(b,c,q)\in\mathcal L_H$ sólo controla la
salida armónica. La pertenencia de las veinticinco coordenadas a
$\operatorname{im}\mathcal D$ exige además
\eqref{x_reciprocidad:img09:relaciones}--\eqref{x_reciprocidad:img09:congruencia}.
Por ejemplo, añadir a $c$ el vector $e_0-e_3$ conserva las tres sumas
orbitales y deja $\Pi_H$ inalterado. Partiendo de $(0,0,0)$ se
obtiene así una terna con $\Pi_H=0$ que incumple
\eqref{x_reciprocidad:img09:relaciones}. Éste es un control negativo exacto de la
compatibilidad del canal, independiente de toda evaluación terminal.

\subsubsection{Determinación y transporte sobre las órbitas del registro}

El Teorema~\ref{x_reciprocidad:img09:imagen} determina $k$ cuando se han producido
sus canales y su carga. Como condición sobre salidas aún no
individualizadas, la misma imagen es un grupo isomorfo a
$\mathbb Z^{12}$. Fijada únicamente una carga $q$, las soluciones
forman una clase afín de
\[
 L_0=\{v\in\mathbb Z^{12}:\sum v_i=0\}
 =\bigoplus_{i=0}^{10}\mathbb Z(e_i-e_{11}).
\]
Por ejemplo, $100\mathbf1+e_0$ y $100\mathbf1+e_1$ son dos
registros distintos, de carga $1201$, con componentes en
$\{0,\ldots,999\}$. Cada uno satisface todas las condiciones de
imagen y todas las congruencias de Hadamard mediante sus propios
canales. Estos testigos comparan las ecuaciones de compatibilidad;
no se les atribuye la condición de historias terminales HMT.

\begin{proposition}[Covariancia cíclica y estabilizador del registro]
\label{x_reciprocidad:img09:covariancia}
Sean $\rho$ la traslación de ventanas en un dominio de registros
$\mathscr R$ y $S$ la traslación anterior en $\mathbb Z^{12}$.
Sobre la órbita de $R\in\mathscr R$, un lector covariante queda
determinado por un vector
\[
 k\in\operatorname{Fix}(S^p),
\]
donde $p\mid12$ es el período mínimo de $R$ bajo $\rho$.
La condición de carga $Q(k)=q$ admite soluciones enteras si y sólo si
$12/p$ divide a $q$; cuando existen, forman una clase afín de rango
$p-1$. En particular, sobre una órbita libre, covariancia y carga
dejan once grados de libertad.
\end{proposition}
\begin{proof}
La regla $F(\rho^jR)=S^jk$ está bien definida exactamente cuando
$S^pk=k$. Esos vectores repiten $12/p$ veces una palabra de
longitud $p$. Su carga es $\frac{12}{p}\sum_{i=0}^{p-1}k_i$.
La divisibilidad y el rango se siguen de esta fórmula. La
covariancia de los canales resulta de $D_aS=SD_a$ y $QS=Q$.
Sobre $\mathcal L_H$, la acción correspondiente es
$\widehat S_H=\mathcal H_{12}S\mathcal H_{12}/4$;
por la Proposición~\ref{x_reciprocidad:img09:triangulo}, ambas reconstrucciones
entrelazan estas acciones.
\end{proof}

Si el dominio lleva un origen marcado, debe utilizarse la acción que
transporta también esa marca. La proposición se refiere a esa acción
cíclica declarada; no presupone la periodicidad de la historia
enriquecida. La naturalidad frente a truncamientos posteriores al corte
$108$ se obtiene, para cualquier lector $F_{108}$ ya definido, mediante
$F_N=F_{108}\circ\tau_{N,108}$. Esta naturalidad conserva la regla
de evaluación inicial; por sí sola no elige una de las reglas posibles
sobre los primeros ciento ocho eventos.

\subsubsection{Evaluación por eventos y composición de contribuciones}

La caracterización anterior proporciona una factorización de la extracción mediante dos
lecturas efectivas del registro:
\begin{equation}
 \omega:\mathscr R\longrightarrow
 \{w\in\mathbb Z^{12}:\sum w_i=0\},\qquad
 q_{\rm reg}:\mathscr R\longrightarrow\mathbb Z,
 \quad q_{\rm reg}(R)+T(\omega(R))\equiv0\pmod {12}.
 \label{x_reciprocidad:img09:contrato-productor}
\end{equation}
Estas aplicaciones deben evaluarse sobre los eventos, la orientación,
la incidencia y la memoria que produjo APP--TRIT--TPK. Una vez dadas sus
acciones sobre esos datos, la composición
\[
 R\longmapsto(w,q)\longmapsto
 ((I+S+S^2)w,(I+S+S^2+S^3)w,q)
 \xrightarrow{\Pi_H}u\xrightarrow{\mathcal H_{12}/4}k
\]
está determinada y satisface todos los controles inversos. La fórmula
$w=c-Sb$ también permite conectar un productor que ya entregue ambos
canales: es una comprobación sobre su salida, anterior a toda comparación
con el testigo terminal.

\begin{proposition}[Acumulación de contribuciones y unicidad prospectiva]
\label{x_reciprocidad:img09:acumulacion}
Supóngase que el estado inicial y cada evento $a_t$ del registro
producen, mediante reglas fijadas previamente, contribuciones
$(\delta w_t,\delta q_t)$ con
\[
 \sum_i\delta w_{t,i}=0,\qquad
 \delta q_t+T(\delta w_t)\equiv0\pmod {12}.
\]
Con condiciones iniciales compatibles $(w^{(0)},q^{(0)})$, la
actualización
\[
 w^{(t+1)}=w^{(t)}+\delta w_t,\qquad
 q^{(t+1)}=q^{(t)}+\delta q_t
\]
produce un único registro en cada etapa. Su incremento es
\[
 \delta k_{t,i}=
 \frac{\delta q_t+T(\delta w_t)}{12}-P_i(\delta w_t).
\]
La construcción respeta la concatenación de registros, y la lectura
de un prefijo permanece igual bajo prolongaciones posteriores.
\end{proposition}
\begin{proof}
Las condiciones de compatibilidad se conservan por suma. La fórmula
\eqref{x_reciprocidad:img09:inversa} es aditiva en $(w,q)$ y produce el
incremento indicado. La inducción fija sucesivamente cada estado.
La suma sobre dos segmentos es la suma concatenada, y un prefijo recibe
exactamente los mismos sumandos antes y después de prolongar la historia.
\end{proof}

La proposición proporciona un criterio suficiente, no una exigencia de
linealidad para todo lector TPK. Si los eventos se transportan por acciones
no triviales, sus contribuciones pueden expresarse en la carta terminal
mediante el cociclo afín correspondiente. La ecuación de cociclo conserva
entonces la misma independencia de la partición de la trayectoria.
En ambos casos, las reglas de evento y el estado inicial constituyen los
datos constructivos que han de provenir del artículo: escogerlos para
reproducir un vector esperado sustituiría esa procedencia por otra tarea.

El resultado de este desarrollo es la condición exacta de determinación
\eqref{x_reciprocidad:img09:contrato-productor}, su inversión demostrada y sus identidades
de transporte. Su alcance es el registro transversal; las restantes
construcciones de HMT conservan sus dominios y sus pruebas propias.

% END OWNER


\subsection{Dos lecturas racionales: \texorpdfstring{\(\kappa_{36}\)}{kappa36} y \texorpdfstring{\(\kappa_{\mathrm{per}}\)}{kappa periódica}}
\label{x_reciprocidad:subsec:k-kappas}

Una vez construido \(K\), la carta en base mil permite dos operaciones
distintas. La primera lee una sola vuelta; la segunda prolonga el registro dodecafásico
periódicamente. Sea \(B=1000\) y definamos su entero de concatenación
\begin{equation}
 N_K=\sum_{m=1}^{12}K_mB^{12-m}
 =234543140729659824621058914794146601.
 \label{x_reciprocidad:eq:k-numerador}
\end{equation}
Los ceros iniciales de un bloque, como el de 058, son parte de su escritura
de tres cifras. No cambian su valor entero, pero sí permiten concatenar
sin ambigüedad los doce bloques.

\begin{definition}[Lecturas finita y periódica del registro dodecafásico]
\label{x_reciprocidad:def:k-kappas}
La lectura de una vuelta es
\[
 \kappa_{36}=\sum_{m=1}^{12}K_mB^{-m}=\frac{N_K}{B^{12}}.
\]
La prolongación periódica es
\[
 K_j^{\mathrm{per}}=K_{1+((j-1)\bmod12)},\qquad
 \kappa_{\mathrm{per}}=\sum_{j\geq1}K_j^{\mathrm{per}}B^{-j}.
\]
\end{definition}

\begin{proposition}[Valor y período exactos del registro dodecafásico prolongado]
\label{x_reciprocidad:prop:k-periodo}
Se cumple
\begin{equation}
 \kappa_{\mathrm{per}}=\frac{N_K}{B^{12}-1},\qquad
 \kappa_{\mathrm{per}}-\kappa_{36}
 =\frac{N_K}{B^{12}(B^{12}-1)}
 =B^{-12}\kappa_{\mathrm{per}}.
 \label{x_reciprocidad:eq:k-diferencia-kappas}
\end{equation}
La palabra infinita del registro dodecafásico tiene período mínimo doce en base mil y
período mínimo treinta y seis en base diez.
\end{proposition}

\begin{proof}
La suma de la primera vuelta es \(N_KB^{-12}\); cada nueva vuelta
multiplica su contribución por \(B^{-12}\). La serie geométrica produce
\(N_K/(B^{12}-1)\), y la resta da la segunda identidad.

Los doce componentes de \(K\) son distintos. Un período propio del ciclo
identificaría dos de ellos; por tanto el período mínimo de bloques es
doce. La expansión decimal es canónica, pues el registro dodecafásico no tiene una cola
de bloques 999. Si su período mínimo decimal fuera \(d\), dividiría
36. Agrupar de tres en tres produciría un período de bloques
\(d/\gcd(d,3)\). Ningún divisor propio de 36 da doce mediante esa
operación. Luego el período decimal mínimo es 36.
\end{proof}

Ambas lecturas son racionales, pero no son el mismo racional. La primera
termina después de doce bloques; la segunda repite la vuelta. Además,
\(0<\kappa_{\mathrm{per}}-\kappa_{36}<B^{-12}\). Su cercanía no
permite sustituir una por otra en una identidad exacta. En particular,
el cierre finito de alfa utiliza \(\kappa_{36}\), mientras que la
suma de la sección prolongada utiliza \(\kappa_{\mathrm{per}}\).


% BEGIN OWNER /Users/ruben/Documents/New project/output/EXTREMA_MEDIA_RAZON_RECIPROCIDAD_ES_EN_20260918/ARTICULO/ES/source/sections/k_reversibilidad.tex
\subsection{Reversibilidad de la constante racional y cambio de fase}

La constante racional \(\kappa_{\mathrm{per}}\) determina íntegramente
el registro \(K\) cuando se conservan la base, la longitud y el origen
de lectura. Su función numérica y su función estructural se relacionan
mediante una inversión exacta.

\begin{proposition}[Recuperación integral y transformación cíclica]
Sea \(I(K)=\sum_{j=1}^{12}K_j1000^{12-j}\). Entonces
\[
 I(K)=(1000^{12}-1)\kappa_{\mathrm{per}}.
\]
La expansión entera de \(I(K)\) en doce bloques de base \(1000\)
recupera todos los \(K_j\), incluidos los ceros iniciales de cada bloque.
Si \(\rho(K)=(K_2,\ldots,K_{12},K_1)\), se cumple
\begin{equation}
 \kappa(\rho K)=1000\kappa(K)-K_1.
 \label{x_reciprocidad:eq:k-rotacion-racional}
\end{equation}
\end{proposition}
\begin{proof}
La primera identidad invierte la definición de la evaluación periódica.
La división euclídea iterada por \(1000\), con doce posiciones fijadas,
recupera el vector. Por otra parte,
\[
 I(\rho K)=1000I(K)-K_1(1000^{12}-1).
\]
Dividir por \(1000^{12}-1\) da \eqref{x_reciprocidad:eq:k-rotacion-racional}.
\end{proof}

La recuperación compone, por tanto, \(\kappa\mapsto K\mapsto U\)
y las transformaciones \(D_3K,D_4K,Q\) ya construidas. La evaluación
racional conserva el registro firmado y las diferencias cíclicas que
dependen de él. La profundidad total y los operadores de transporte de
una historia permanecen en sus coordenadas propias: la inversión anterior
tiene por dominio el registro integral dodecafásico.

Escribamos \(\kappa_j=\kappa(\rho^jK)\), con índices módulo \(12\).
La misma identidad da
\[
 K_{j+1}=1000\kappa_j-\kappa_{j+1},\qquad
 \sum_{j=0}^{11}\kappa_j=\frac{\sum_{j=1}^{12}K_j}{999}
 =\frac{6263}{999}.
\]
La fase actúa de manera exacta sobre la constante: su valor es fijo
en la lectura canónica, mientras sus rotaciones siguen una ley afín.
La suma de la órbita es invariante bajo la elección del origen.

También la incidencia posterior dispone de una expresión en estas
coordenadas. El umbral cúbico determina
\[
 B_K=\{j+1:1000\kappa_j-\kappa_{j+1}\ge729\}
 =\{4,6,9,10\}.
\]
Se recupera la tétrada a partir de la órbita racional porque ésta
recupera primero las componentes enteras. El soporte firmado asociado
a \(\alpha\) utiliza además sus registros regionales y la normalización
de frontera. La relación entre coordenada e incidencia queda así
expresada mediante operaciones identificables.

Esta es una definición operacional de la función de \(K\): proporciona
una coordenatización integral reversible del registro orientado, cuya
evaluación periódica y cuyas funciones de incidencia participan en
construcciones posteriores. El valor, la transformación y el dominio
se conservan conjuntamente.

% END OWNER


\subsection{Periodicidad y conúcleo del operador de transporte}
\label{x_reciprocidad:subsec:k-clase-acarreo}

La periodicidad que acabamos de demostrar pertenece al registro dodecafásico. No implica
periodicidad de las coordenadas regionales, de los acarreos de la clausura
ni de alfa. Es útil precisar esta separación mediante un segundo
cociente, puramente integral.

Sea \(S\) el desplazamiento cíclico de doce coordenadas, con orientación
fijada, y sea \(b\geq2\) una base entera. El operador de acarreo
periódico es \(\partial_b=S-bI\). Si una identidad periódica se escribe
\[
 A=P+E-\Phi-K+\partial_bC,
\]
entonces las transformaciones
\[
 K\longmapsto K+\partial_bh,\qquad C\longmapsto C+h
\]
conservan su lado derecho. Esta invariancia compara representantes de una
clase de acarreo; no es una autorización para modificar el registro dodecafásico ya
seleccionado por el registro.

\begin{proposition}[Cociente del acarreo periódico]
\label{x_reciprocidad:prop:k-cociente-acarreo}
Con la orientación anterior,
\[
 \operatorname{coker}(S-bI)\cong\mathbb Z/(b^{12}-1)\mathbb Z.
\]
\end{proposition}

\begin{proof}
El módulo cíclico se presenta como
\(\mathbb Z[t]/(t^{12}-1)\). Imponer \(S=bI\) equivale a imponer
\(t=b\), y deja la relación \(b^{12}-1=0\). La evaluación de los
coeficientes en \(b\), con el orden de concatenación elegido, representa
la clase resultante.
\end{proof}

El denominador \(B^{12}-1\) de \(\kappa_{\mathrm{per}}\) y este
cociente proceden de la misma longitud cíclica, pero sus objetos siguen
siendo distintos: uno es una evaluación racional y el otro un cociente
de un módulo de acarreos. La clausura finita de alfa será una cadena
orientada con frontera derecha. Su prolongación transportará una
frontera remota compatible, no un acarreo que se identifique por decreto
entre la primera y la decimotercera posición.

Queda así preparado el capítulo siguiente. El registro dodecafásico se ha fijado antes
de resolver la recurrencia de alfa; su reconstrucción por Hadamard es
entera y única; sus diferencias transversales recuperan el registro; y
sus dos evaluaciones racionales tienen dominios y usos separados. Ninguno
de estos hechos utiliza alfa como entrada, y ninguno decide todavía el
tipo aritmético de la salida de la clausura completa.
